% ============================================================
%  BrushHeroes — Manual de Usuario
%  Pontificia Universidad Javeriana · Ingeniería de Sistemas
% ============================================================
\documentclass[12pt, a4paper]{article}

% ── Paquetes ─────────────────────────────────────────────────
\usepackage[spanish]{babel}
\usepackage[utf8]{inputenc}
\usepackage[T1]{fontenc}
\usepackage{lmodern}
\usepackage{geometry}
\usepackage{graphicx}
\usepackage{float}

% Imagen segura: muestra un recuadro gris si el archivo aún no existe.
% Cuando tenga las capturas en img/, este comando las mostrará automáticamente.
\let\realincludegraphics\includegraphics
\renewcommand{\includegraphics}[2][]{%
  \IfFileExists{#2}%
    {\realincludegraphics[#1]{#2}}%
    {\fbox{\begin{minipage}[c][3cm][c]{0.5\textwidth}%
       \centering\small\color{gray}%
       \textbf{[imagen pendiente]}\\[4pt]\texttt{#2}%
     \end{minipage}}}%
}
\usepackage{caption}
\usepackage{subcaption}
\usepackage{xcolor}
\usepackage{titlesec}
\usepackage{fancyhdr}
\usepackage{hyperref}
\usepackage{enumitem}
\usepackage{tcolorbox}
\usepackage{booktabs}
\usepackage{array}
\usepackage{parskip}
\usepackage{setspace}
\usepackage{microtype}

\geometry{
  top=2.5cm, bottom=2.5cm,
  left=3cm,  right=2.5cm
}

% ── Paleta de colores ─────────────────────────────────────────
\definecolor{heroBlue}{HTML}{2563EB}
\definecolor{heroGreen}{HTML}{22C55E}
\definecolor{heroOrange}{HTML}{F97316}
\definecolor{heroRed}{HTML}{EF4444}
\definecolor{heroGray}{HTML}{64748B}
\definecolor{heroBg}{HTML}{F0F9FF}
\definecolor{heroLightGray}{HTML}{F8FAFC}

% ── Tipografía de secciones ───────────────────────────────────
\titleformat{\section}
  {\Large\bfseries\color{heroBlue}}
  {\thesection.}{0.8em}{}[\vspace{-4pt}\rule{\linewidth}{0.6pt}]

\titleformat{\subsection}
  {\large\bfseries\color{heroBlue!80!black}}
  {\thesubsection.}{0.6em}{}

\titleformat{\subsubsection}
  {\normalsize\bfseries\color{heroGray}}
  {\thesubsubsection.}{0.5em}{}

\titlespacing{\section}{0pt}{18pt}{8pt}
\titlespacing{\subsection}{0pt}{12pt}{4pt}
\titlespacing{\subsubsection}{0pt}{8pt}{2pt}

% ── Cajas de aviso ────────────────────────────────────────────
\tcbuselibrary{skins, breakable}

\newtcolorbox{notebox}[1][Nota]{
  colback=heroBg, colframe=heroBlue,
  fonttitle=\bfseries\small, title=#1,
  left=6pt, right=6pt, top=4pt, bottom=4pt,
  arc=4pt, boxrule=0.8pt
}

\newtcolorbox{tipbox}[1][Consejo]{
  colback=heroGreen!8, colframe=heroGreen!55!black,
  fonttitle=\bfseries\small, title=#1,
  left=6pt, right=6pt, top=4pt, bottom=4pt,
  arc=4pt, boxrule=0.8pt
}

\newtcolorbox{warnbox}[1][Importante]{
  colback=heroOrange!8, colframe=heroOrange,
  fonttitle=\bfseries\small, title=#1,
  left=6pt, right=6pt, top=4pt, bottom=4pt,
  arc=4pt, boxrule=0.8pt
}

% ── Encabezado y pie de página ────────────────────────────────
\pagestyle{fancy}
\fancyhf{}
\fancyhead[L]{\small\color{heroGray}BrushHeroes}
\fancyhead[R]{\small\color{heroGray}\leftmark}
\fancyfoot[C]{\small\color{heroGray}\thepage}
\renewcommand{\headrulewidth}{0.4pt}
\renewcommand{\headrule}{\hbox to\headwidth{\color{heroBlue!30}\leaders\hrule height \headrulewidth\hfill}}

% ── Estilo de figuras ─────────────────────────────────────────
\captionsetup{
  font=small,
  labelfont={bf,color=heroGray},
  textfont={color=heroGray},
  belowskip=4pt
}

% ── Metadatos PDF ─────────────────────────────────────────────
\hypersetup{
  colorlinks=true,
  linkcolor=heroBlue,
  urlcolor=heroBlue,
  pdftitle={BrushHeroes - Manual de Usuario},
  pdfauthor={Equipo BrushHeroes - PUJ},
  pdfsubject={Manual de Usuario}
}

% ── Comando para placeholder de imagen ───────────────────────
% Útil mientras no se tienen las capturas finales.
% Uso: \imgplaceholder{ancho}{alto}{nombre_archivo.png}{descripción}
\newcommand{\imgplaceholder}[4]{%
  \begin{figure}[H]
    \centering
    \fbox{\begin{minipage}[c][#2][c]{#1}
      \centering
      \small\color{heroGray}
      \textbf{IMAGEN PENDIENTE}\\[4pt]
      \texttt{img/#3}\\[4pt]
      #4
    \end{minipage}}
  \end{figure}%
}

% ─────────────────────────────────────────────────────────────
\begin{document}

% ── PORTADA ──────────────────────────────────────────────────
\begin{titlepage}
  \centering
  \vspace*{1.5cm}

  {\small\color{heroGray} PONTIFICIA UNIVERSIDAD JAVERIANA}\\[0.2cm]
  {\small\color{heroGray} FACULTAD DE INGENIERÍA — INGENIERÍA DE SISTEMAS}\\[0.6cm]
  \rule{0.7\textwidth}{0.8pt}

  \vspace{1.8cm}

  % Reemplace el bloque siguiente por \includegraphics cuando tenga el logo:
  % \includegraphics[width=0.42\textwidth]{img/portada_logo.png}
  \begin{tcolorbox}[
    colback=heroBlue!8, colframe=heroBlue,
    width=0.55\textwidth, halign=center,
    arc=8pt, boxrule=1.2pt,
    top=18pt, bottom=18pt
  ]
    \centering
    {\LARGE\bfseries\color{heroBlue} BrushHeroes}\\[4pt]
    {\small\color{heroGray} Videojuego educativo de higiene bucal}
  \end{tcolorbox}
  % NOTA: reemplazar el tcolorbox anterior por el logo cuando esté disponible.

  \vspace{1.8cm}

  {\huge\bfseries Manual de Usuario}\\[0.6cm]
  {\large\color{heroGray} Guía completa para el jugador}

  \vspace{1.5cm}
  \rule{0.5\textwidth}{0.4pt}\\[0.8cm]

  \begin{tabular}{r l}
    \color{heroGray}Plataforma:   & Android (Realidad Aumentada) \\
    \color{heroGray}Versión:      & 1.0 \\
    \color{heroGray}Fecha:        & \today \\
  \end{tabular}

  \vfill
  {\small\color{heroGray} Proyecto de Grado · Ingeniería de Sistemas\\
  Pontificia Universidad Javeriana · \the\year}
\end{titlepage}

% ── TABLA DE CONTENIDOS ───────────────────────────────────────
\renewcommand{\contentsname}{Tabla de contenidos}
\tableofcontents
\newpage

% =============================================================
\section{Introducción}
% =============================================================

\textbf{BrushHeroes} es una aplicación móvil educativa diseñada para motivar a los niños a
desarrollar y mantener buenos hábitos de higiene bucal. A través de minijuegos interactivos y un
sistema de progreso gamificado —inspirado en aplicaciones como Duolingo— el jugador aprende la
técnica correcta de cepillado y uso del hilo dental de una manera entretenida y accesible.

\begin{figure}[H]
  \centering
  \includegraphics[width=0.50\textwidth]{img/intro_menu_general.png}
  % IMAGEN: captura completa del menú principal (pestaña Minijuegos, scroll al tope).
  \caption{Pantalla principal de BrushHeroes.}
  \label{fig:intro_menu}
\end{figure}

\subsection{¿A quién está dirigido?}

BrushHeroes está diseñado principalmente para niños entre 4 y 10 años, aunque puede ser
utilizado por cualquier persona. Para los más pequeños se recomienda la supervisión de un adulto
durante la configuración inicial y las primeras sesiones de juego.

\subsection{Minijuegos disponibles}

\begin{itemize}[leftmargin=1.8cm, itemsep=4pt]
  \item \textbf{Cepillado con Realidad Aumentada (AR):} usa la cámara frontal del dispositivo
        para superponer bacterias y un cepillo virtual sobre el rostro del jugador, guiándolo por
        las seis zonas de cepillado de la boca.
  \item \textbf{Cepillado interactivo:} práctica de la técnica de cepillado mediante un
        minijuego en orientación horizontal.
  \item \textbf{Hilo dental:} aprende a usar el hilo dental correctamente con un minijuego
        interactivo.
\end{itemize}

\subsection{Sistema de progreso}

Además de los minijuegos, BrushHeroes incluye:

\begin{itemize}[leftmargin=1.8cm, itemsep=4pt]
  \item \textbf{Racha de días:} registra cuántos días consecutivos ha jugado el usuario.
  \item \textbf{Estrellas:} recompensas obtenidas al completar sesiones.
  \item \textbf{Misiones diarias:} sesión de mañana y sesión de tarde.
  \item \textbf{Calendario de actividad:} visualización del historial del mes con código de color.
\end{itemize}

\begin{notebox}
  Esta aplicación \textbf{no reemplaza} la instrucción de un profesional de salud bucal.
  Su propósito es reforzar hábitos de higiene de forma lúdica.
\end{notebox}

% =============================================================
\section{Requisitos del sistema}
% =============================================================

Para instalar y ejecutar BrushHeroes correctamente, el dispositivo debe cumplir con los
siguientes requisitos mínimos:

\begin{table}[H]
  \centering
  \caption{Requisitos mínimos del dispositivo.}
  \label{tab:requisitos}
  \renewcommand{\arraystretch}{1.3}
  \begin{tabular}{>{\bfseries\color{heroGray}}p{4.2cm} p{8.5cm}}
    \toprule
    \color{black}Componente & \color{black}Requisito \\
    \midrule
    Sistema operativo  & Android 9.0 (API 28) o superior \\
    Procesador         & Octa-core 2.0\,GHz o superior \\
    Memoria RAM        & 3\,GB mínimo (4\,GB recomendado) \\
    Almacenamiento     & 300\,MB de espacio libre \\
    Cámara frontal     & Requerida para el minijuego AR (mínimo 5\,MP) \\
    Soporte AR         & Compatible con Google ARCore \\
    Conexión a red     & No requerida para jugar \\
    \bottomrule
  \end{tabular}
\end{table}

\begin{warnbox}
  El \textbf{minijuego de Realidad Aumentada} requiere que el dispositivo sea compatible con
  \textbf{Google ARCore}. Puede consultar la lista oficial de dispositivos compatibles en:
  \url{https://developers.google.com/ar/devices}
\end{warnbox}

% =============================================================
\section{Instalación}
% =============================================================

\subsection{Instalación desde archivo APK}

BrushHeroes se distribuye como un archivo \texttt{.apk} que debe instalarse manualmente en
el dispositivo Android. Siga estos pasos:

\begin{enumerate}[leftmargin=1.8cm, itemsep=6pt]
  \item Descargue el archivo \texttt{BrushHeroes.apk} desde el enlace proporcionado por su
        institución educativa y guárdelo en el dispositivo.

  \item En su dispositivo, vaya a \textbf{Ajustes $\rightarrow$ Seguridad} (o
        \textbf{Ajustes $\rightarrow$ Aplicaciones $\rightarrow$ Instalar apps desconocidas}) y habilite la opción
        \textit{«Permitir instalación de fuentes desconocidas»}.

        \begin{figure}[H]
          \centering
          \includegraphics[width=0.38\textwidth]{img/install_fuentes_desconocidas.png}
          % IMAGEN: ajustes de Android mostrando la opción de fuentes desconocidas habilitada.
          \caption{Activar instalación desde fuentes desconocidas.}
          \label{fig:install_fuentes}
        \end{figure}

  \item Abra el archivo \texttt{BrushHeroes.apk} desde el administrador de archivos del
        dispositivo. Aparecerá un diálogo de confirmación; toque \textbf{Instalar}.

        \begin{figure}[H]
          \centering
          \includegraphics[width=0.38\textwidth]{img/install_dialogo.png}
          % IMAGEN: diálogo de Android "¿Desea instalar esta aplicación?" con botón Instalar.
          \caption{Diálogo de instalación del APK.}
          \label{fig:install_dialogo}
        \end{figure}

  \item Espere a que finalice la instalación. Al completarse, toque \textbf{Abrir} para
        iniciar la aplicación, o búsquela en su lanzador de aplicaciones bajo el nombre
        \textbf{BrushHeroes}.
\end{enumerate}

\begin{tipbox}
  Si el dispositivo muestra una advertencia sobre aplicaciones de fuentes desconocidas,
  es normal para aplicaciones distribuidas fuera de la Play Store. El archivo es seguro.
\end{tipbox}

\subsection{Primer inicio}

Al abrir la aplicación por primera vez, el sistema generará automáticamente un perfil de
jugador con datos de ejemplo para el calendario de actividad. No es necesario crear una cuenta.

% =============================================================
\section{Menú principal}
% =============================================================

Al iniciar la aplicación se mostrará el menú principal. La pantalla está organizada en dos
pestañas accesibles desde la barra de navegación inferior: \textbf{Minijuegos} y
\textbf{Calendario}.

\begin{figure}[H]
  \centering
  \includegraphics[width=0.50\textwidth]{img/menu_completo.png}
  % IMAGEN: captura completa del menú — pestaña Minijuegos visible, scroll al tope,
  %         mascota, racha, estrellas y las 3 tarjetas todas visibles.
  \caption{Menú principal — pestaña Minijuegos.}
  \label{fig:menu_completo}
\end{figure}

\subsection{Panel del héroe}

En la parte superior de la pestaña Minijuegos se encuentra el \textbf{panel del héroe}, que
muestra de un vistazo el estado de progreso del jugador:

\begin{figure}[H]
  \centering
  \includegraphics[width=0.52\textwidth]{img/menu_hero_panel.png}
  % IMAGEN: primer plano del panel superior — mascota animada a la izquierda,
  %         llama de racha con número de días, estrella con contador,
  %         chips "Mañana" y "Tarde" (idealmente uno con ✓ y otro sin él),
  %         mensaje motivacional abajo.
  \caption{Panel del héroe con racha, estrellas y misiones diarias.}
  \label{fig:hero_panel}
\end{figure}

\begin{itemize}[leftmargin=1.8cm, itemsep=4pt]
  \item \textbf{Mascota:} personaje animado cuyo estado refleja el desempeño del jugador.
        Si la racha es alta, la mascota celebra; si no ha jugado en varios días, se muestra triste.
  \item \textbf{Racha}: número de días consecutivos en los que el jugador ha
        completado al menos una sesión. ¡Intente mantenerla activa!
  \item \textbf{Estrellas}: total de estrellas acumuladas completando sesiones de juego.
  \item \textbf{Misión Mañana / Tarde:} indica si ya completó la sesión de cepillado de la
        mañana o de la tarde en el día actual. Aparece una marca $\checkmark$ al completarla.
  \item \textbf{Mensaje motivacional:} texto dinámico que cambia según la racha actual.
\end{itemize}

\subsection{Tarjetas de minijuegos}

Debajo del panel del héroe se encuentran las tres tarjetas de acceso a los minijuegos.
Cada tarjeta muestra el nombre del juego y una breve descripción.

\begin{figure}[H]
  \centering
  \includegraphics[width=0.52\textwidth]{img/menu_cards.png}
  % IMAGEN: las 3 tarjetas visibles — ARCard, BrushCard y FlossCard,
  %         haciendo scroll down si es necesario para verlas todas.
  \caption{Tarjetas de los tres minijuegos disponibles.}
  \label{fig:menu_cards}
\end{figure}

Para ingresar a un minijuego, toque el botón \textbf{«Jugar»} en la tarjeta correspondiente.

\begin{notebox}
  El minijuego de \textbf{Cepillado AR} requiere la cámara frontal del dispositivo.
  Los minijuegos de \textbf{Cepillado interactivo} y \textbf{Hilo Dental} se ejecutan en
  orientación horizontal; el dispositivo rotará automáticamente al ingresar.
\end{notebox}

\subsection{Pestaña Calendario}

Toque el ícono \textbf{Calendario} en la barra inferior para ver el historial de actividad
del mes actual.

\begin{figure}[H]
  \centering
  \includegraphics[width=0.52\textwidth]{img/menu_calendario.png}
  % IMAGEN: pestaña Calendario completa — encabezado con mes/año y flechas de navegación,
  %         grid con días coloreados, leyenda de colores en la parte inferior.
  \caption{Pestaña Calendario con historial de actividad.}
  \label{fig:calendario}
\end{figure}

Cada celda del calendario representa un día del mes y se colorea según el desempeño:

\begin{table}[H]
  \centering
  \caption{Código de colores del calendario.}
  \label{tab:colores_calendario}
  \renewcommand{\arraystretch}{1.3}
  \begin{tabular}{>{\centering\arraybackslash}p{1.5cm} p{9cm}}
    \toprule
    \textbf{Color} & \textbf{Significado} \\
    \midrule
    \textcolor{heroGreen}{\rule{1.2cm}{0.4cm}}  & $\geq$ 80\% de sesiones completadas (excelente) \\
    \textcolor{yellow!60!orange}{\rule{1.2cm}{0.4cm}} & $\geq$ 50\% de sesiones completadas (bien) \\
    \textcolor{heroOrange}{\rule{1.2cm}{0.4cm}} & Al menos una sesión completada \\
    \textcolor{heroRed!70}{\rule{1.2cm}{0.4cm}} & Día sin actividad (perdido) \\
    \textcolor{heroBlue}{\rule{1.2cm}{0.4cm}}   & Día de hoy \\
    \textcolor{gray!30}{\rule{1.2cm}{0.4cm}}    & Días futuros (sin datos) \\
    \bottomrule
  \end{tabular}
\end{table}

Use las flechas \textbf{«\textless»} y \textbf{«\textgreater»} en el encabezado del calendario
para navegar entre meses anteriores y el mes actual.

% =============================================================
\section{Minijuego AR: Cepillado con Realidad Aumentada}
% =============================================================

El minijuego de Realidad Aumentada es la experiencia principal de BrushHeroes. Usa la cámara
frontal del dispositivo para detectar el rostro del jugador y superponer elementos virtuales
(bacterias y un cepillo de dientes) sobre su imagen en tiempo real.

\subsection{Cómo acceder}

Desde el menú principal, toque la tarjeta \textbf{«Cepillado AR»} y luego el botón
\textbf{«Jugar»}. La aplicación abrirá la cámara frontal y mostrará la pantalla de inicio del
minijuego.

\subsection{Detección de rostro}

Al entrar al minijuego, la aplicación buscará automáticamente su rostro. Mientras no lo
detecte, los botones de inicio permanecerán bloqueados y se mostrará un aviso en pantalla.

\begin{figure}[H]
  \centering
  \begin{subfigure}[b]{0.42\textwidth}
    \includegraphics[width=\textwidth]{img/ar_buscando_rostro.png}
    % IMAGEN: pantalla de inicio AR con overlay visible "Buscando rostro..."
    %         (botones de modo en gris/deshabilitados).
    \caption{Buscando rostro\ldots}
  \end{subfigure}
  \hfill
  \begin{subfigure}[b]{0.42\textwidth}
    \includegraphics[width=\textwidth]{img/ar_rostro_detectado.png}
    % IMAGEN: misma pantalla pero con cara detectada y overlay ocultado
    %         (botones de modo en color/habilitados).
    \caption{Rostro detectado.}
  \end{subfigure}
  \caption{Estados de la pantalla de inicio: buscando y encontrando el rostro.}
  \label{fig:ar_deteccion}
\end{figure}

\begin{tipbox}
  Para una detección óptima: colóquese en un lugar bien iluminado, a unos 30--50\,cm de
  la pantalla, y asegúrese de que su cara esté completamente visible y centrada.
\end{tipbox}

\subsection{Selección de modo de juego}

Una vez detectado el rostro, podrá elegir entre dos modos de juego:

\begin{figure}[H]
  \centering
  \includegraphics[width=0.50\textwidth]{img/ar_seleccion_modo.png}
  % IMAGEN: pantalla de inicio con los dos botones activos y visibles:
  %         "Dinámico" y "Solo guía".
  \caption{Pantalla de selección de modo.}
  \label{fig:ar_modos}
\end{figure}

\begin{itemize}[leftmargin=1.8cm, itemsep=6pt]
  \item \textbf{Modo Dinámico:} el jugador controla el cepillo virtual arrastrando su dedo
        por la pantalla. Debe limpiar todas las zonas de la boca antes de que se agoten los
        3 minutos. Recomendado para niños que ya conocen la dinámica del juego.
  \item \textbf{Modo Solo Guía:} el cepillo se mueve automáticamente sobre cada zona,
        mostrando la técnica correcta de cepillado. No tiene límite de tiempo. Ideal para
        el primer contacto con el juego o para niños más pequeños.
\end{itemize}

\subsection{Modo Dinámico}

\subsubsection{Objetivo}

Limpie las \textbf{seis zonas de la boca} eliminando todas las bacterias de cada zona antes
de que se acaben los \textbf{3 minutos} (180 segundos).

\subsubsection{Cómo jugar}

\begin{figure}[H]
  \centering
  \includegraphics[width=0.52\textwidth]{img/ar_dinamico_gameplay.png}
  % IMAGEN: gameplay modo Dinámico — cara del jugador con bacterias animadas visibles
  %         sobre alguna zona, cepillo virtual posicionado sobre ellas,
  %         HUD en la parte superior (puntos, 6 dots de zona, cronómetro).
  \caption{Modo Dinámico en acción: bacterias sobre la zona activa y cepillo virtual.}
  \label{fig:ar_dinamico}
\end{figure}

\begin{enumerate}[leftmargin=1.8cm, itemsep=5pt]
  \item Al comenzar la sesión, las bacterias aparecerán sobre la \textbf{primera zona activa}
        de su boca.
  \item \textbf{Arrastre su dedo por la pantalla} para mover el cepillo virtual. El cepillo
        sigue el movimiento de su dedo en tiempo real.
  \item Pase el cepillo sobre las bacterias para eliminarlas. Escuchará un sonido al
        destruir cada bacteria.
  \item Cuando todas las bacterias de la zona activa sean eliminadas, se avanzará
        automáticamente a la siguiente zona. Una notificación indicará el cambio.
  \item Complete las \textbf{seis zonas} para ganar la sesión.
\end{enumerate}

\subsubsection{Las seis zonas de cepillado}

El juego cubre las mismas zonas que recomienda limpiar un odontólogo:

\begin{table}[H]
  \centering
  \caption{Zonas de cepillado del minijuego AR.}
  \label{tab:zonas}
  \renewcommand{\arraystretch}{1.3}
  \begin{tabular}{>{\bfseries}p{4cm} p{8cm}}
    \toprule
    Zona & Ubicación \\
    \midrule
    Superior izquierda  & Molares y premolares superiores del lado izquierdo \\
    Superior derecha    & Molares y premolares superiores del lado derecho \\
    Frontal superior    & Incisivos y caninos superiores (dientes frontales arriba) \\
    Inferior izquierda  & Molares y premolares inferiores del lado izquierdo \\
    Inferior derecha    & Molares y premolares inferiores del lado derecho \\
    Frontal inferior    & Incisivos y caninos inferiores (dientes frontales abajo) \\
    \bottomrule
  \end{tabular}
\end{table}

\subsection{Modo Solo Guía}

\begin{figure}[H]
  \centering
  \includegraphics[width=0.52\textwidth]{img/ar_guiado_gameplay.png}
  % IMAGEN: modo Solo Guía — cepillo animado moviéndose con un movimiento de vaivén
  %         sobre una zona de la boca. Sin HUD de puntos ni cronómetro visibles.
  \caption{Modo Solo Guía: el cepillo se mueve automáticamente sobre cada zona.}
  \label{fig:ar_guiado}
\end{figure}

En este modo el jugador no necesita tocar la pantalla. El cepillo virtual recorre cada zona
de la boca con el movimiento correcto de cepillado. Las bacterias se eliminan solas de forma
secuencial, permitiendo al jugador observar la técnica.

\begin{itemize}[leftmargin=1.8cm, itemsep=4pt]
  \item No hay cronómetro ni puntuación en este modo.
  \item La sesión avanza zona por zona de manera automática.
  \item Una guía visual en la parte inferior indica la zona que se está cepillando en ese momento.
\end{itemize}

\begin{tipbox}
  Use el \textbf{Modo Solo Guía} la primera vez para aprender el recorrido correcto de
  cepillado antes de intentar el Modo Dinámico.
\end{tipbox}

\subsection{HUD durante el juego}

Durante una sesión en \textbf{Modo Dinámico}, la barra superior de la pantalla muestra la
información de la sesión en tiempo real:

\begin{figure}[H]
  \centering
  \includegraphics[width=0.72\textwidth]{img/ar_hud_detalle.png}
  % IMAGEN: primer plano del HUD superior — panel de puntuación a la izquierda,
  %         6 puntos de zona en el centro (algunos verdes, uno amarillo, resto gris),
  %         cronómetro a la derecha mostrando segundos restantes.
  \caption{HUD de juego: puntuación, progreso de zonas y cronómetro.}
  \label{fig:ar_hud}
\end{figure}

\begin{itemize}[leftmargin=1.8cm, itemsep=4pt]
  \item \textbf{Puntuación} (izquierda): puntos acumulados en la sesión actual.
        Aumenta al eliminar bacterias.
  \item \textbf{Indicadores de zona} (centro): seis puntos que representan las seis zonas.
        \begin{itemize}
          \item \textcolor{heroGreen}{\textbf{Verde}}: zona completada.
          \item \textcolor{yellow!70!orange}{\textbf{Amarillo}}: zona activa (en progreso).
          \item \textcolor{gray}{\textbf{Gris}}: zona pendiente.
        \end{itemize}
  \item \textbf{Cronómetro} (derecha): segundos restantes de la sesión. Solo visible en
        Modo Dinámico.
\end{itemize}

\subsection{Guía de zona}

En la parte inferior de la pantalla, durante toda la sesión, se muestra una animación que
indica la zona que debe cepillarse en ese momento.

\subsection{Notificación de zona}

Al completar una zona, aparece brevemente una notificación en pantalla que muestra el nombre
de la siguiente zona a cepillar.

\begin{figure}[H]
  \centering
  \includegraphics[width=0.50\textwidth]{img/ar_zona_notificacion.png}
  % IMAGEN: overlay de notificación de zona — fondo semitransparente oscuro,
  %         nombre de la nueva zona en texto grande, icono/imagen de la zona.
  \caption{Notificación al completar una zona y pasar a la siguiente.}
  \label{fig:ar_zona_notif}
\end{figure}

\subsection{Pausa}

Durante el juego puede pausar la sesión tocando el botón de pausa. El tiempo se detendrá
y aparecerá el menú de pausa con las siguientes opciones:

\begin{figure}[H]
  \centering
  \includegraphics[width=0.50\textwidth]{img/ar_pausa.png}
  % IMAGEN: pantalla de pausa — overlay oscuro con opciones
  %         "Continuar" y "Salir al Menú".
  \caption{Menú de pausa.}
  \label{fig:ar_pausa}
\end{figure}

\begin{itemize}[leftmargin=1.8cm, itemsep=4pt]
  \item \textbf{Continuar:} reanuda la sesión desde donde se pausó.
  \item \textbf{Salir al Menú:} abandona la sesión y regresa al menú principal.
\end{itemize}

\subsection{Pantalla de resultado}

Al finalizar la sesión (por completar todas las zonas o por agotar el tiempo), se muestra
la pantalla de resultado:

\begin{figure}[H]
  \centering
  \begin{subfigure}[b]{0.42\textwidth}
    \includegraphics[width=\textwidth]{img/ar_resultado_exito.png}
    % IMAGEN: pantalla de resultado exitoso — "¡Excelente!"
    %         "¡Limpiaste todos tus dientes!" + puntuación final + animación celebración.
    \caption{Resultado exitoso.}
  \end{subfigure}
  \hfill
  \begin{subfigure}[b]{0.42\textwidth}
    \includegraphics[width=\textwidth]{img/ar_resultado_fallo.png}
    % IMAGEN: pantalla de resultado fallido — "¡Tiempo agotado!"
    %         "¡Inténtalo de nuevo, tú puedes!" + puntuación final.
    \caption{Tiempo agotado.}
  \end{subfigure}
  \caption{Pantallas de resultado al finalizar la sesión AR.}
  \label{fig:ar_resultado}
\end{figure}

Desde esta pantalla puede:
\begin{itemize}[leftmargin=1.8cm, itemsep=4pt]
  \item \textbf{Reintentar:} iniciar una nueva sesión del mismo minijuego.
  \item \textbf{Volver al menú:} regresar al menú principal.
\end{itemize}

% =============================================================
\section{Minijuego: Cepillado interactivo}
% =============================================================

El minijuego de Cepillado interactivo se ejecuta en orientación horizontal. Al tocarlo desde
el menú, el dispositivo rotará automáticamente.

\begin{figure}[H]
  \centering
  \includegraphics[width=0.72\textwidth]{img/brush_gameplay.png}
  % IMAGEN: gameplay del minijuego BrushingGame en landscape —
  %         mostrar la pantalla principal del juego con los elementos interactivos.
  \caption{Minijuego de Cepillado interactivo.}
  \label{fig:brush_game}
\end{figure}

\subsection{Objetivo}

% TODO: completar con la descripción específica del minijuego de cepillado.
% Descripción provisional:
Siga las instrucciones en pantalla para cepillar correctamente cada zona de los dientes.
Complete la rutina completa de cepillado para ganar estrellas.

\subsection{Cómo volver al menú}

Utilice el botón de retroceso del dispositivo o el botón de menú disponible en pantalla
para regresar al menú principal.

% =============================================================
\section{Minijuego: Hilo Dental}
% =============================================================

El minijuego de Hilo Dental también se ejecuta en orientación horizontal.

\begin{figure}[H]
  \centering
  \includegraphics[width=0.72\textwidth]{img/floss_gameplay.png}
  % IMAGEN: gameplay del minijuego DentalFlossGame en landscape —
  %         mostrar la pantalla principal con los elementos del juego.
  \caption{Minijuego de Hilo Dental.}
  \label{fig:floss_game}
\end{figure}

\subsection{Objetivo}

% TODO: completar con la descripción específica del minijuego de hilo dental.
% Descripción provisional:
Aprenda a pasar el hilo dental correctamente entre los dientes siguiendo las indicaciones
visuales en pantalla. Complete todas las piezas dentales para finalizar la sesión.

\subsection{Cómo volver al menú}

Utilice el botón de retroceso del dispositivo o el botón de menú disponible en pantalla
para regresar al menú principal.

% =============================================================
\section{Sistema de progreso}
% =============================================================

BrushHeroes lleva un registro del desempeño del jugador a lo largo del tiempo para motivar
la constancia en la higiene bucal.

\subsection{Racha de días}

\begin{figure}[H]
  \centering
  \includegraphics[width=0.44\textwidth]{img/progreso_racha.png}
  % IMAGEN: primer plano del badge de racha en el panel del héroe —
  %         llama animada con el número de días (idealmente 5+ para que se vea bien).
  \caption{Indicador de racha de días consecutivos.}
  \label{fig:racha}
\end{figure}

La \textbf{racha} aumenta en uno cada día en que el jugador completa al menos una sesión de
juego. Si se pierde un día sin jugar, la racha vuelve a cero.

La mascota del panel principal refleja el estado de la racha:

\begin{table}[H]
  \centering
  \caption{Estado de la mascota según la racha.}
  \label{tab:mascota}
  \renewcommand{\arraystretch}{1.3}
  \begin{tabular}{>{\bfseries}p{4cm} p{8cm}}
    \toprule
    Racha & Estado de la mascota \\
    \midrule
    7 días o más  & Celebrando — animación de festejo \\
    3 a 6 días    & Feliz — animación de saludo \\
    1 a 2 días    & Neutral \\
    0 días        & Triste \\
    \bottomrule
  \end{tabular}
\end{table}

\subsection{Estrellas}

Las \textbf{estrellas} se acumulan completando sesiones de juego. El total se muestra en el
panel del héroe junto al ícono de estrella.

\subsection{Misiones diarias}

Cada día el jugador tiene dos misiones:
\begin{itemize}[leftmargin=1.8cm, itemsep=4pt]
  \item \textbf{Misión de Mañana:} completar una sesión de juego durante la mañana.
  \item \textbf{Misión de Tarde:} completar una sesión de juego durante la tarde.
\end{itemize}

Al completar una misión, el chip correspondiente en el panel del héroe mostrará una marca $\checkmark$.

\subsection{Calendario de actividad}

El \textbf{calendario} de la pestaña Calendario muestra el historial del mes actual con
código de colores (ver tabla~\ref{tab:colores_calendario}). Use las flechas de navegación
para ver meses anteriores.

\begin{notebox}
  Los datos del calendario se guardan localmente en el dispositivo. Si desinstala la
  aplicación, el historial se perderá.
\end{notebox}

% =============================================================
\section{Preguntas frecuentes}
% =============================================================

\subsubsection*{¿Por qué la cámara no detecta mi rostro?}
Asegúrese de estar en un lugar con buena iluminación y de que su cara esté completamente
visible frente a la cámara, a unos 30--50\,cm de distancia. Si el problema persiste,
verifique que su dispositivo sea compatible con Google ARCore
(\url{https://developers.google.com/ar/devices}).

\subsubsection*{¿Qué hago si el juego se ve lento o con interrupciones?}
Cierre otras aplicaciones en segundo plano y reinicie la aplicación. Si el dispositivo no
cumple los requisitos mínimos (ver sección~\ref{tab:requisitos}), el rendimiento puede verse afectado.

\subsubsection*{¿Se pierde mi progreso si cierro la aplicación?}
No. El progreso (racha, estrellas, historial) se guarda automáticamente en el dispositivo
cada vez que regresa al menú principal.

\subsubsection*{¿Se pierde mi progreso si desinstalo la app?}
Sí. Los datos se almacenan localmente. Desinstalar la aplicación eliminará todo el
historial y el progreso acumulado.

\subsubsection*{¿Necesito conexión a internet para jugar?}
No. BrushHeroes funciona completamente sin conexión a internet una vez instalado.

\subsubsection*{¿Puedo usar auriculares?}
Sí. El juego incluye efectos de sonido y música de fondo. Los auriculares son opcionales.

\subsubsection*{¿El juego funciona con luz tenue?}
La detección de rostro para el minijuego AR requiere iluminación razonable. En condiciones
de poca luz el reconocimiento facial puede fallar o ser intermitente.

\subsubsection*{¿Cuánto tiempo dura una sesión?}
En \textbf{Modo Dinámico} la sesión dura máximo 3 minutos (180 segundos).
En \textbf{Modo Solo Guía} no hay límite de tiempo.

\subsubsection*{¿Para qué sirve el Modo Solo Guía?}
Está diseñado para enseñar la técnica de cepillado correcta de manera pasiva. Es ideal
para niños que usan la aplicación por primera vez o para adultos que la utilizan como
material de demostración.

% ── FIN DEL DOCUMENTO ────────────────────────────────────────
\vfill
\begin{center}
  \small\color{heroGray}\rule{0.5\textwidth}{0.4pt}\\[6pt]
  BrushHeroes · Manual de Usuario · Versión 1.0\\
  Pontificia Universidad Javeriana · Ingeniería de Sistemas · \the\year
\end{center}

\end{document}
